% [原文 p.287]

设 $C_Z\subset\PP^{n+1}$ 是以 $(0:\cdots:0:1)$ 为顶点的 $Z$ 上的锥。射影 $\pi\colon(x_0:\cdots:x_{n+1})\mapsto(x_n:x_{n+1})$ 定义态射 $f_Z\colon X_Z\to\PP^1$。其中 $X_Z:=B_DC_Z$ 是 $C_Z$ 沿 $D$ 的爆破。若 $s\ne0$ 则 $f_Z^{-1}(s:t)\cong Z$，若 $s=0$ 则 $f_Z^{-1}(s:t)\cong C_D$。

把这个构造应用于 Veronese 曲面 $V\subset\PP^5$ 与正规 4 次有理直纹曲面 $(F_2\cong)S\subset\PP^5$。两者的任意一般超平面截面都是正规 4 次有理曲线 $D\subset\PP^4$。可以构成满足 $f_V^{-1}(0:1)\cong C_D$ 且 $f_V^{-1}(1:t)\cong V$（或 $f_S^{-1}(0:1)\cong C_D$ 且 $f_S^{-1}(1:t)\cong S$）的曲面族 $f_V\colon X_V\to\PP^1$（或 $f_S\colon X_S\to\PP^1$）。

$C_D$ 是正规 4 次有理曲线上的锥，其极小奇点消解是直纹曲面 $F_4$。在顶点附近，锥 $C_D$ 与商奇点 $\C^2/\mu_4(1,1)$ 同型，特别是对数终端（lt）的。

注意自交点数 $(K_V^2)=9$、$(K_S^2)=8$、$(K_{C_D}^2)=9$。第一个族中，纤维的标准类的自交点数是常数，但第二个族在原点处发生跳跃。由此可见，仅凭纤维不能充分描述族的好坏。

$X_V$ 只有标准奇点（唯一的奇点是 Veronese 曲面上的锥的顶点，是终端商奇点 $\C^3/\mu_2(1,1,1)$），但若注意到 $K_{C_S}$ 连 Q-Cartier 都不是，就能稍许理解问题的性质。\hfill$\square$

由这个例子可知，对于非奇异曲线上的曲线族 $f\colon X\to C$，仅假设各纤维是 lc 的并不充分。为保证 $X$ 是标准的，需要追加条件。本质上由定理 5.50，只需假设 $K_{X/C}$ 是 Q-Cartier 即可。

对于一般的底 $f\colon X\to Y$，这一条件变得更加微妙。只要 $Y$ 是非奇异的，假设 $K_{X/Y}$ 是 Q-Cartier 就足够了。一般情形，特别是底 $Y$ 不是约化的时候，相应的条件变得技术性（参看 [Kol94b]）。\hfill$\square$

利用半稳定 MMP，至少就曲面而言，任意 1-参数族都可以紧化。

% [原文 p.288]
\begin{theoremp}{7.62}
设 $C$ 是 $\C$ 上的非奇异曲线，$C^0\subset C$ 是开子概形。设 $f^0\colon(X^0,B^0)\to C^0$ 是使 $K_{X^0}+B^0$ 为 $f^0$-丰富的固有 lc 态射。

此时，可以适当地取非奇异曲线 $C'$、有限满射 $h\colon C'\to C$，以及作为 $(X^0,B^0)\times_C C'\to h^{-1}(C^0)$ 的延拓的固有 lc 态射 $f'\colon(X',B')\to C'$，使 $K_{X'}+B'$ 成为 $f'$-丰富的。用图式表示则如下。
\[
\begin{tikzcd}[column sep=2.1em,row sep=1.3em]
(X^0,B^0) & (X^0,B^0)\times_C C' \arrow[l] & (X',B') \arrow[l,hook]\\
C^0 \arrow[u,"f^0"] & C^0\times_C C' \arrow[l] \arrow[u] & C' \arrow[l,hook] \arrow[u,"f'"']
\end{tikzcd}
\]
\end{theoremp}

\begin{proof}
把 $f^0\colon(X^0,B^0)\to C^0$ 延拓为射影态射 $f\colon(X,B)\to C$。应用半稳定约化定理 7.19，得到使 $(X^0,B^0)\times_C C'$ 具有半稳定奇点消解的 $h\colon C'\to C$。在此应用定理 7.10，便得到 $f'\colon(X',B')\to C'$。
\end{proof}

% [Tutor Note: 定理 7.62 表明必须允许基变换：只靠原来的底 $C$ 往往不能完成延拓，半稳定约化定理 7.19 正是为此提供了有限满射 $C'\to C$。]

\par\medskip\noindent\textbf{7.63}（商奇点的形变空间）\quad 设 $(0\in H)$ 是奇点。如定理 4.60 所述，$(0\in H)$ 具有除去局部解析同构后唯一确定的半普遍形变空间 $\operatorname{Def}(H)$。

定理 4.61 对超曲面奇点的情形给出具体构造。

然而一般情形下，记述半普遍形变空间非常困难。下面说明如何用半稳定 MMP 计算 2 维商奇点 $H$ 的半普遍形变空间 $\operatorname{Def}(H)$ 的不可约分支的个数。即使在这种情形也非常微妙。先从辅助结果开始。\hfill$\square$

\begin{lemmap}{7.64}
设 $0\in X$ 是 3 维终端奇点，$0\in S\subset X$ 是 Cartier 除子。此时下述等价。
\begin{enumerate}[label=(\arabic*),leftmargin=2.4em,itemsep=1pt,topsep=2pt,parsep=0pt]
\item $(X,S)$ 是 plt 的。
\item $0\in S$ 是商奇点。
\item $0\in S$ 是 Du Val 奇点，或解析同构于适当的 $(xy+z^{dn}=0)/\mu_n(1,-1,r)$（$(r,n)=1$）。
\end{enumerate}
\end{lemmap}

\begin{proof}
(1) 与 (2) 等价由定理 5.50 与命题 4.18 可知。由 (3) 得 (2) 是明显的。逆用定理 5.43 便不难。

% [原文 p.289]
详细情况请参看 [KSB88, 3.10]。
\end{proof}

\begin{theorem*}{7.65}
（[KSB88, 3.9]）设 $(0\in H)$ 是曲面奇点。则下述两个集合之间存在一一对应。
\begin{enumerate}[label=(\arabic*),leftmargin=2.4em,itemsep=1pt,topsep=2pt,parsep=0pt]
\item $\operatorname{Def}(H)$ 的不可约分支。
\item 固有双有理态射 $f\colon H'\to H$（除同型外），使 $K_{H'}$ 为 $f$-丰富，且 $H'$ 的各奇点都是引理 7.64(3) 的奇点之一。
\end{enumerate}
\end{theorem*}

% [Tutor Note: 定理 7.65 的「一一对应」把形变论（$\operatorname{Def}(H)$ 的分支）与双有理几何（好的 $H'$）连起来：分支越少，$H$ 的极小消解越刚性。]

依据 [KSB88]，集合（定理 7.65(2)）的计算方法另有不同。[BC94, Ste93] 进一步计算这一集合，开发出能求得关于 $\operatorname{Def}(H)$ 的许多信息的非常有效的算法。把这些想法进一步一般化，就能对非常多的曲面奇点记述 $\operatorname{Def}(H)$ 的不可约分支（参看 [Kol91b, Sec.6]）。

\textbf{[证明概要]}\quad 这里只说明如何从不可约分支 $Z\subset\operatorname{Def}(H)$ 构成双有理态射 $f\colon H'\to H$。为此需要下述结果：$\operatorname{Def}(H)$ 的各不可约分支都包含光滑化（smoothing）（参看 [Art74]）。准确叙述该结果，就是存在平坦形变
\[
\begin{array}{ccc}
H & \subset & Y\\
\downarrow & & \downarrow{g}\\
0 & \in & C
\end{array}
\]
且 $g\colon Y\setminus H\to C$ 光滑，对应态射 $C\to\operatorname{Def}(H)$ 的像含于 $Z$。

根据需要作基变换，便可假设 $g\colon Y\to C$ 具有半稳定奇点消解。由定理 7.9 与定理 7.10，$Y$ 具有极小模型与标准模型。把引理 7.66 应用于极小模型，便知在标准模型 $f\colon X\to Y$ 上 $H':=f^*H$ 满足定理 7.65(2) 的条件。

还有几点需要确认：$H'$ 不依赖于 $g\colon Y\to C$ 的选择，以及所有 $H'$ 都能这样得到。这些可归结为 2 维有理奇点与其部分奇点消解（partial resolution）两者的形变性质。详见 [KSB88, Sec.3]。\hfill$\blacksquare$
% [原文 p.290]
\begin{lemmap}{7.66}
设 $Y$ 是正规簇，$H\subset Y$ 是 Cartier 除子，$f\colon X\to Y$ 是固有双有理态射。假定下述条件。
\begin{enumerate}[label=(\arabic*),leftmargin=2.4em,itemsep=1pt,topsep=2pt,parsep=0pt]
\item $(H,0)$ 是正规且 klt 的。
\item $X$ 是 Q-分解的，且 $K_X$ 是 $f$-nef 的。
\item $(X,f^*H)$ 是 dlt 的。
\end{enumerate}
此时，$f^*H=f_*^{-1}H$，且 $(X,f^*H)$ 在 $f^*H$ 附近是 plt 的。
\end{lemmap}

\begin{proof}
用 $f$-例外除子 $E$ 写成 $f^*H=f_*^{-1}H+E$。设 $y\in H$ 是余维 1 的点。此时，$H$ 以及 $Y$ 在 $y$ 处非奇异。由于 $K_X$ 是 $f$-nef 的，由此可知 $f$ 在 $y$ 上同型（引理 3.39）。特别地，$(f_*^{-1}H)\cap E$ 是 $f|_{f_*^{-1}H}$ 例外的。

由于 $f_*^{-1}H$ 是余维 2 处的 Cartier 除子（推论 5.56），可得次数关系
\[
K_{f_*^{-1}H}\equiv(K_X+f_*^{-1}H)|_{f_*^{-1}H}\equiv f(K_X-E)|_{f_*^{-1}H}.
\]
另一方面，令 $K_{f_*^{-1}H}\equiv f^*K_H+F$，则 $F$ 是 $f$-例外的，且 $|F|$ 是有效的。

组合这些结果，得到 $F+E|_{f_*^{-1}H}\equiv_f K_X|_{f_*^{-1}H}$，由此特别可知 $F+E|_{f_*^{-1}H}$ 是 $f$-nef 的。由引理 3.39，这意味着 $-F-E|_{f_*^{-1}H}$ 是有效的。$-F$ 的各不可约分支的系数 $<1$，且若 $E\ne0$，则 $E|_{f_*^{-1}H}$ 的至少一个不可约分支的系数 $\ge1$。因此得到 $E=0$、$F\le0$、$f^*H=f_*^{-1}H$。

由 $K_{f_*^{-1}H}-F\equiv f^*K_H$ 知 $(f_*^{-1}H,-F)$ 是 klt 的。由于 $-F$ 是有效的，$(f_*^{-1}H,0)$ 是 klt 的（推论 2.35）。这样由定理 5.50，$(X,f^*H)$ 在 $f^*H$ 附近是 plt 的。
\end{proof}

\par\medskip\noindent\textbf{7.67}（曲面族的同时奇点消解）\quad 在 §4.3 中已述：Du Val 奇点的任意平坦族，至少作有限满射的基变换后，就具有同时奇点消解。对于曲面的任意平坦族 $g\colon Y\to S$，找出保证同时奇点消解存在的条件是有趣的问题。这个问题本质上局部于 $Y$，但以固有曲面族为对象时能得到最明快的答案。这一判别法可以局部化，但需要注意（参看 [KSB88, 2.25]）。\hfill$\square$

\begin{theoremp}{7.68}
（[KSB88]）设 $S$ 是 $\C$ 上的概形，$g\colon Y\to S$ 是固有平坦态射，其任意纤维都是只有孤立奇点的曲面。对各点 $s\in S$，设 $\bar Y_s\to Y_s$ 是纤维的极小奇点消解。此时下述成立。
\begin{enumerate}[label=(\arabic*),leftmargin=2.4em,itemsep=1pt,topsep=2pt,parsep=0pt]
\item $s\mapsto(K_{Y_s}^2)$ 是下半连续的。
\item $g$ 经有限满射 $S'\to S$ 的基变换后具有同时奇点消解，当且仅当 $s\mapsto(K_{Y_s}^2)$ 是局部常值函数。
\end{enumerate}
\end{theoremp}

% [原文 p.291]
\textbf{[证明概要]}\quad 只对 $S=C$ 是非奇异曲线的情形证明。由此容易导出定理 7.68(1)。

一旦关于同时奇点消解的几个一般结果（参看 [KSB88, Sec.2]）得到确立，便知定理 7.68(2) 是其归结。

作基变换 $C'\to C$ 后，半稳定奇点消解存在。为简化记号，假设 $g\colon Y\to C$ 自身具有半稳定奇点消解。对此应用定理 7.9，得到
\[
X\xrightarrow{f}Y\xrightarrow{g}C,
\]
其中 $X\to C$ 是 dlt 的，$K_X$ 是 $f$-nef 的。用 $X_c$ 表示 $c\in C$ 上 $g\circ f$ 的纤维。$X_c$ 对泛点 $c\in C$ 是 $Y_c$ 的极小奇点消解，但对特定的值（例如 $0\in C$）纤维可能可约或奇异。写成 $X_0=X_0'+E_0$。其中 $X_0'$ 是 $Y_0$ 的双有理变换，$E$ 是 $f$-例外除子。此时成立下述次数关系。
\[
\begin{aligned}
(K_{\bar Y_c}^2)&=(K_{X_c}^2)=(K_{X_0}^2)=((K_{X_0}|_{X_0'})^2)+((K_{X_0}|_E)^2)\\
&\ge((K_{X_0}|_{X_0'})^2)=((K_{X_0'}+E|_{X_0'})^2)\ge(K_{\bar Y_0}^2).
\end{aligned}
\]
其中最后的不等式由引理 7.69。这就证明了 (1)。等号成立的必要充分条件是 $E=0$，且 $X_0$ 至多只有 Du Val 奇点；这一情形下再作基变换则同时奇点消解存在（定理 4.28）。\hfill$\blacksquare$

\begin{lemmap}{7.69}
设 $S$ 是固有曲面，$\bar S\to S$ 是极小奇点消解。$f_1\colon S_1\to S$ 是固有双有理态射，假定 $(S_1,0)$ 是 klt 的。$\Delta_1$ 是有效 $f_1$-例外的

% [原文 p.292]
Q-除子，且 $K_{S_1}+\Delta_1$ 是 $f_1$-nef 的。此时，$((K_{S_1}+\Delta_1)^2)\ge(K_{\bar S}^2)$ 成立。并且，等号成立的必要充分条件是 $\Delta_1=0$ 且 $S_1$ 至多只有 Du Val 奇点。
\end{lemmap}

\begin{proof}
设 $g_1\colon S_1\to S_2$ 是 $(S_1,0)$ 的 $S$ 上极小模型。特别地，令 $K_{S_1}\equiv g_1^*K_{S_2}+\Delta_2$，则 $\Delta_2$ 是有效的。由此 $K_{S_1}+\Delta_1\equiv g_1^*K_{S_2}+(\Delta_1+\Delta_2)$。因此对 $S_1\to S$ 应用引理 7.70，得到 $((K_{S_1}+\Delta_1)^2)\ge(K_{S_2}^2)$。

设 $g_3\colon S_3\to S_2$ 是极小奇点消解。此时令 $K_{S_3}\equiv g_3^*K_{S_2}-\Delta_3$，则 $\Delta_3$ 是有效的（推论 4.3）。特别地，由引理 7.70 得到 $(K_{S_2}^2)\ge(K_{S_3}^2)$。

最后，设 $g_4\colon S_3\to\bar S$ 是 $S_3$ 的 $S$ 上极小模型。与上面同样得到 $(K_{S_3}^2)\ge(K_{\bar S}^2)$。

由以上可知，等号成立的必要充分条件是 $\Delta_1=0$、$S_2=S_1$ 且 $S_1$ 只有 Du Val 奇点，并且 $S_3=\bar S$。
\end{proof}

\begin{lemmap}{7.70}
设 $g\colon U\to V$ 是正规固有曲面之间的双有理态射。设 $N$ 是 $g$-nef 的 Q-除子，$E$ 是有效 $g$-例外的 Q-除子，且 $N+E$ 是 $g$-nef 的。此时，$((N+E)^2)\ge(N^2)$。并且，等号成立当且仅当 $E=0$。
\end{lemmap}

\begin{proof}
$((N+E)^2)=(N^2)+(N\cdot E)+(N+E\cdot E)\ge(N^2)$ 成立。若等号成立，则 $(N\cdot E)=(N+E\cdot E)=0$，从而特别地 $(E^2)=0$。由引理 3.40，这意味着 $E=0$。
\end{proof}

\section{相关结果概述}
\takeaway{本节不再给出证明，而是以文献表的形式概观与半稳定 MMP 相关的结果：翻转、GIT、小平维数、丰饶性、藤田猜想，以及 Fano、Calabi--Yau、椭圆簇等方向。}

本节的目的是概述相关结果与本书所述方向的发展。比起精确叙述结果并作详细解说，更以给出文献的相当完整的表为目的。

\par\medskip\noindent\textbf{7.71}（翻转）\quad 主要结果之一，正如已多次提到的，是 3 维翻转的存在定理。无边界情形的终端 3 维簇的情形由 [Mor88] 给出。[KM92] 发展这一方法，对只有终端奇点的 3 维翻转进行了分类。
% [原文 p.293]
这一结果有几个应用。例如，射影非奇异 3 维簇的多重亏格（plurigenus）的形变不变性，以及一般型 3 维簇的双有理等价类的粗模空间（coarse moduli space）的存在等。（参看 [KM92, Chap.12]、7.82。）

关于 3 维 lc 对 $(X,\Delta)$ 的更一般的翻转，由 [Sho92] 处理。该论文的许多想法由研讨班笔记 [K+92, Chaps.16--22] 进一步发展。此外，该笔记通过把 3 维 lc 翻转的存在归结为无边界终端情形 [Mor88]，给出其另一证明 [K+92, Chaps.4--8]。归结的方法与 §6.4 所用的讨论类似。

4 维以上翻转的存在尚不清楚。这大概是极小模型理论中最重要的未解决问题。已确立的结果中最有意义的大概是 toric 簇的情形 [Rei83a]，以及其他特殊代数群作用的簇的情形 [BK94]。\hfill$\square$

% [Tutor Note: 7.71 指出四维以上翻转仍未知，这是极小模型理论最重要的未解决问题；本章的三维结果为它提供了方法原型。]

\par\medskip\noindent\textbf{7.72}（翻转与几何不变量论（GIT））\quad 自翻转被发现以来，人们已理解翻转在许多状况下出现。其中最重要的是，翻转在 [Mum65] 的几何不变量论中自然出现。群对概形的几何商空间，只有选取称为极化的丰富除子之后才能定义。选取不同的极化就得到不同的商空间。这些商空间彼此双有理等价，其双有理映射往往由翻转或逆翻转的复合得到。一般理论由 [Tha96, DH98, MW97] 发展。\hfill$\square$

\par\medskip\noindent\textbf{7.73}（小平维数）\quad 设 $X$ 是非奇异射影簇。若 $m$ 充分大且充分可除，则 $|mK_X|\colon X\dashrightarrow\PP$ 的像的双有理等价类不依赖于 $m$ 的选择。这个像的维数称为 $X$ 的\textbf{小平维数}（Kodaira dimension），用 $\kappa(X)$ 表示。小平维数、特别是纤维空间情形下其行为的研究，由饭高茂开始。这一方向的结果大多在极小模型理论被研究之前。[Uen75] 是良好的解说书，[Mor87] 是其后续结果的概说。\hfill$\square$

\par\medskip\noindent\textbf{7.74}（丰度（abundance）与标准环）\quad 已在 §3.2 与 §6.1 中触及标准环的有限生成性与丰度。原型是下述有限生成性问题。

非奇异射影簇的标准环 $\bigoplus_{m=0}^{\infty}H^0(X,\cO_X(mK_X))$ 是否有限生成？

关于 3 维簇的最初重要结果，是在 $\kappa(X)\le2$ 的假定下证明有限生成性的 [Fuj86]。把 3 维问题归结为 2 维 lc 对 $(S,\Delta)$ 的这一方法，是 lc 对理论初期的成功之一。

如 §6.1 所述，若假定适当的 MMP，则丰度意味着有限生成性（不过逆命题大概不成立）。关于 3 维簇的丰度，$\kappa=1$ 与 $\kappa=2$ 的情形需要不同手法等，已判明极其困难。证明由川又与宫冈的一系列论文确立。[K+92, Chaps.9--15] 介绍简化的讨论并给出文献表。关于 3 维 lc 对 $(X,\Delta)$ 的丰度由 [KMM94a] 证明。

% [原文 p.294]
还有一个有趣的问题如下。给出满足下述性质的常数 $c(n)$。若 $X$ 是 $n$ 维非奇异射影簇且满足 $\kappa(X)\ge0$，则对适当的 $0<m\le c(n)$，$|mK_X|\ne\emptyset$ 成立。

关于 3 维簇，$\kappa=0$ 的情形由 [Kaw86] 处理。$\kappa=1$ 的情形是森（未发表），$\kappa=2$ 的情形由 [Kol94a] 给出。$\kappa=3$ 的情形，尽管已有若干部分结果 [Kol86a, Luo94]，仍未解决。\hfill$\square$

% [Tutor Note: 注意区分「丰度」（abundance）与「丰富」（ample）：前者面对的是 $K_X$ 的丰度猜想，后者是除子的正性，中译按术语表分别处理。]

\par\medskip\noindent\textbf{7.75}（藤田猜想（Fujita's conjecture））\quad 设 $X$ 是非奇异射影簇，$L$ 是 $X$ 上的丰富除子。藤田隆夫猜想：若 $m\ge\dim X+1$，则 $|K_X+mL|$ 是自由的。更一般地，可以提出 $K_X$ 与 $L$ 的哪种线性组合是 nef、自由、丰富、非常丰富的问题。这一方向已得到许多结果，其中 Reider、Demailly、Ein--Lazarsfeld、川又雄二郎、Angehrn--Siu、Helmke、Smith 等人的贡献重要。[Laz96] 的讲义是非常好的解说。

目前最强的结果是：藤田猜想在维数 $\le4$ 时成立 [Kaw97]。任意维数下，若 $m>\frac12\dim X(\dim X+1)$，则 $|K_X+mL|$ 是自由的 [AS95]。[Kol97, secs.5--6] 有最近结果的概说。
% [原文 p.295]
这些结果所发展的手法也可应用于 Abel 簇的研究。其最初是 [Kol95, Chap.17]，而最强的结果由 [EL97] 给出。为了传达这些应用的氛围，叙述其中一个结果。

设 $(A,\Theta)$ 是不可分解的主极化 Abel 簇。此时 $(A,\Theta)$ 是标准的。换言之，$\Theta$ 是正规的，且至多只有有理奇点。\hfill$\square$

\par\medskip\noindent\textbf{7.76}（Fano 纤维空间的双有理几何）\quad 在 §3.8 中说明了标准模型唯一、极小模型在余维 1 意义下唯一。与此相对，若从 $X$ 出发的 MMP 以 Fano 纤维空间终结，则选择不同的极射线往往会以完全不同的 Fano 纤维空间终结。判定两个 Fano 纤维空间是否彼此双有理等价似乎非常困难。

对的 MMP 可以从两条途径对解决这个问题作出贡献。首先，若适当选取边界 $\Delta$，则对 $(X,\Delta)$ 作 MMP 得到的 Fano 纤维空间有时会变得特别简单。这一方法由 [Ale91, Cor96] 探求。[CF93] 把这一手法应用于其他情形。

其次，把 MMP 的手法与 Sarkisov 的想法组合起来，[Cor94] 开发出把 Fano 纤维空间之间的双有理映射分解为初等双有理映射的方法。不过，理解这样的双有理映射并不容易，只在少数情形得到完全解决。\hfill$\square$

\par\medskip\noindent\textbf{7.77}（3 维 Fano 簇）\quad MMP 促使人们详细研究其过程中产生的簇。这样的簇的一个类是终端 3 维 Fano 簇。非奇异 3 维 Fano 簇已经分类（参看 [Isk80, MM81]）。它由 104 个不同的形变型组成。关于带奇点的 3 维 Fano 簇，所知甚少。$\rho=1$ 的 Q-分解终端 3 维 Fano 簇至多只有有限个形变型 [Kaw92a]。同样的结果对标准 3 维 Fano 簇也成立，且无需 Q-分解性与 $\rho=1$ 的假定 [KMM94b]。目前尚不知道制作分类表的可行方法。[Fle89] 计算的例子表明表会变得庞大。\hfill$\square$

\par\medskip\noindent\textbf{7.78}（3 维卡拉比--丘簇）\quad 卡拉比--丘簇是标准类数值
% [原文 p.296]
平凡的正规簇。（不同的人对奇点的限制等有所不同。）这样的簇在没有终端奇点时是小平维数 0 的极小模型，因此构成重要的类。（若假定丰度猜想，逆也成立。）然而，近年来对卡拉比--丘簇的兴趣几乎都源于它与弦论的关系（例如参看 [Mor97] 及其广泛文献表）。这些簇的形变型的数量庞大，甚至可能无限。[Rei87a] 的猜想是试图把所有 3 维卡拉比--丘簇相互联系起来的尝试。这一方向的重要结果是 [NS95] 确立的许多 3 维卡拉比--丘簇的光滑化可能性（smoothability）。[SBW94] 给出与 3 维 Abel 簇相关的 3 维卡拉比--丘簇的特征。\hfill$\square$

\par\medskip\noindent\textbf{7.79}（3 维椭圆簇）\quad 当簇 $X$ 具有使一般纤维为椭圆曲线的态射 $f\colon X\to Z$ 时，称 $X$ 是\textbf{椭圆的}（elliptic）。若 $X$ 是椭圆的，则 $\kappa(X)<\dim X$ 成立；反之，$\kappa(X)=\dim X-1$ 的 $X$ 双有理等价于椭圆簇。正如椭圆曲面的小平标准丛公式所显示的，理解奇异纤维的行为似乎很重要。关于 3 维，[Nak91, DG94] 给出了几乎完全的记述。\hfill$\square$

\par\medskip\noindent\textbf{7.80}（复簇的 MMP）\quad 对固有簇或紧复簇定义 MMP 存在若干困难。基本问题大概是：曲线锥作为对象本来就不正确。这起因于有效曲线可以数值平凡。

然而，所有这些情形或许都能用几何学上明确的方法记述，并存在对应的收缩态射。[Kol91a] 讨论了一些特殊情形。

关于 Kähler 簇，这些问题不产生，MMP 的相当类似物似乎可以起作用。[CP97] 在这一方向给出了重要结果。\hfill$\square$

\par\medskip\noindent\textbf{7.81}（正标数或混合标数下的 MMP）\quad 人们相信 MMP 在正标数或混合标数下也起作用。

% [原文 p.297]
[Kol91a] 开始了正标数下 3 维簇的研究。[Kee98] 考察了在正标数下处理无基点性的手段，并在许多情形证明了收缩定理。关于翻转几乎未知。

[Kaw94] 对曲面的半稳定族证明了正标数或混合标数情形下的全部 MMP。\hfill$\square$

\par\medskip\noindent\textbf{7.82}（多重亏格的形变不变性）\quad 最近，[Siu98] 证明了：对一般型簇的光滑族 $f\colon X\to S$，纤维的多重亏格 $h^0(X_s,\cO_{X_s}(mK_{X_s}))$ 对任意 $m\ge1$ 都是 $S$ 上的局部常值函数。这一结果由 [Kaw98] 推广到纤维具有标准奇点的情形，由 [Nak98] 推广到纤维约化的情形。此外，[Kaw98] 证明了标准奇点的小形变是标准的，[Nak98] 证明了终端奇点的小形变是终端的。\hfill$\square$

% [Tutor Note: 7.6 节不给出证明，其价值在于文献导航；读时可把每条结果与第 3、6 章的方法对照，不必逐句推敲。]

\chapterreview{第 7 章把第 3 章建立的锥定理、翻转与极小模型纲领推广到带边界的半稳定情形，是全书从三维终端奇点的极小模型理论走向高维族与退化问题的枢纽。全章的主线是：把「沿纤维的奇点行为」打包成射的性质（lc 射、dlt 射），再在此基础上重跑一遍极小模型纲领。这条路线既是第 3 章锥定理的自然延伸，也解释了为何第 7 章的技术难度明显上升。

第 7.1 节给出 lc 射、dlt 射的定义，并叙述半稳定 MMP 定理：dlt 边界全程保持，终点或为 $f^m$-nef，或为相对 Fano 纤维空间。第 7.2 节的半稳定约化定理把任意族经分歧覆盖化为半稳定族；这是全部讨论的起点，也使「三维半稳定翻转」成为可逐例验证的对象。第 7.3、7.4 节处理翻转：特殊半稳定翻转与一般半稳定翻转的存在分别由定理 7.32、7.40 给出，终止性由定理 7.7 保证；一般半稳定翻转则采用「对数奇点消解 + 修正边界」的策略，核心是定理 7.51 及其证明中的主张 7.53——追加的约化边界吸收掉所有非特殊翻转，再用定理 7.44 把边界减回去。标准模型的存在（定理 7.51）与唯一性（定理 3.52）配合相对无基点定理，给出 MMP 的终点。

第 7.5 节把上述机器用于三个具体问题：稳定曲线模空间的紧化（7.56 与稳定曲线族）、曲面商奇点半普遍形变空间的不可约分支（定理 7.65 的一一对应），以及曲面族的同时奇点消解判别法（定理 7.68：$(K^2)$ 的局部常值性）。例 7.61 用两个曲面族说明 $K^2$ 会在退化纤维处跳跃，正呼应定理 7.68。第 7.6 节不给出证明，以文献表的形式概观翻转、GIT、小平维数、丰度、藤田猜想、Fano 簇、卡拉比--丘簇与椭圆簇等方向的进展。

读本章时宜与第 3 章对照：三维无边界情形是「终端奇点 + 极射线收缩」的干净故事，带边界后，边界成分与纤维退化使技术细节急剧增加。三个容易忽略的辨析点是：尾符三档（证明紧跟、证明在后文、书中不证）、$\operatorname{Def}(H)$ 的分支与好的 $H'$ 的一一对应、以及定理 7.62 与 7.68 中基变换不可省。

评注：定理 7.51 的证明是全书技术高峰之一，不必强求一次读懂；7.6 节的文献表则可当作延伸阅读的地图，按兴趣逐条追溯。若把本章与第 6 章的翻转构造对读，能更清楚地看到「半稳定」假设究竟换来了什么。对初学者，建议先掌握定理 7.9 与 7.10 的陈述，再回头细读定理 7.51 的证明结构。}